\documentclass[11pt]{article}
\usepackage[T1]{fontenc}
\usepackage{lmodern,amsmath,amssymb,amsthm,mathtools}
\usepackage[margin=28mm]{geometry}
\usepackage{microtype,booktabs,array}
\usepackage[hidelinks]{hyperref}
\setlength{\emergencystretch}{3em}
\allowdisplaybreaks[2]
\newtheorem{theorem}{Theorem}[section]
\newtheorem{lemma}[theorem]{Lemma}
\newtheorem{proposition}[theorem]{Proposition}
\newtheorem{corollary}[theorem]{Corollary}
\newtheorem{remark}[theorem]{Remark}
\newcommand{\R}{\mathbb R}
\newcommand{\Z}{\mathbb Z}
\newcommand{\E}{\mathcal E}
\newcommand{\tr}{\operatorname{tr}}
\newcommand{\rank}{\operatorname{rank}}
\newcommand{\TV}{\operatorname{TV}}
\newcommand{\mean}{\operatorname{mean}}
\newcommand{\Id}{\mathrm I}
\newcommand{\normn}[1]{\|#1\|_{V}}
\newcommand{\Prob}{\mathbb P_X}
\newcommand{\Expect}{\mathbb E_X}
\title{Rectangular observations for finite matrix gradients:\\critical budgets and the weak endpoint}
\author{Lluis Eriksson}
\date{10 November 2026 --- manuscript 226B, preparation version}
\hypersetup{pdftitle={Rectangular observations for finite matrix gradients: critical budgets and the weak endpoint},pdfauthor={Lluis Eriksson}}
\begin{document}
\maketitle
\begin{abstract}
For a finite periodic lattice and arbitrary finite symmetric matrices acting on
its actual forward gradient stack, we construct a partition into integer
rectangles whose scaled component means restore coercivity. The output
dimension is $O_D(h(1+F_c(A)))$ for $h$ real components and capped local
eigenchannel budget $F_c(A)$. The product order $hF_c(A)$ is necessary in
worst-case joint growth for common partitions retaining all component means. The quadratic penalty is finite and may
depend on the coefficient profile, including its negative spectral depth.
The proof combines positive spectral replacement, an arbitrary-positive-scalar
balanced-ring construction, and singleton refinement at every affected start
and forward endpoint. We give the full scalar BV, weighted-gap and compressed-tree
proof in an appendix. For positive scalar coefficients with exact mean $\mu$,
a weak inverse-$D/2$ distribution bound gives a logarithmic upper bound on
observation dimension. An explicit family of actual flat coexact fields
forces the same worst-case logarithmic order for every centered differentiable
observation with a finite locally quadratic upper penalty. The mechanisms are
classical; historical priority and non-equivalence to the closest predecessors
remain unestablished. All results concern specified finite component-gradient
forms.
\end{abstract}

\section{Objects, principal results and scope}\label{sec:main}
A small number of spatial averages can restore coercivity even when local
coefficients are very small or have negative channels. Two issues must be
kept separate: how many averages are needed, and how large a penalty must be
placed on them. We control the first by a capped coefficient budget while
allowing the second to depend on the entire finite profile. At the scalar
weak endpoint a logarithmic loss in the coefficient scale is unavoidable in
the worst case.

\subsection{The actual finite form and observation rows}
Throughout, $D\ge3$, $n\ge1$ and $h\ge1$ are integers, $c>0$,
\[
 X=(\Z/n\Z)^D,\qquad V=n^D,\qquad P=D/2,\qquad r=Dh.
\]
We represent $X$ by the fixed chart $\{0,\ldots,n-1\}^D$.
For $u:X\to\R^h$ the forward stack is
\[
 du(x)=\bigl(u(x+e_i)-u(x)\bigr)_{i=1}^D\in\R^r.
\]
All additions in this expression are periodic. Given a finite real symmetric
$r\times r$ matrix $A(x)$ at each original start $x$, define
\begin{equation}\label{eq:actual}
 Q_A(u)=\sum_{x\in X}du(x)^{\mathsf T}A(x)du(x),
 \qquad \normn{u}^2=V^{-1}\sum_x\|u(x)\|^2.
\end{equation}
No common eigenbasis or sign assumption is imposed on $A(x)$.
For a scalar profile $a>0$ we write
\[
 Q_a(u)=\sum_{x,i}a(x)\|u(x+e_i)-u(x)\|^2.
\]
In particular, the coefficient is always evaluated at the forward start;
it is never moved to the endpoint or averaged over an edge.

An integer rectangle is a nonempty product of consecutive intervals in the
fixed chart. It need not be a periodic rectangle crossing the chart boundary.
For such a rectangle $B$ define the vector row
\begin{equation}\label{eq:rows}
 q_B(u)=V^{-1/2}|B|^{-1/2}\sum_{x\in B}u(x)
       =\sqrt{|B|/V}\,\mean_B u.
\end{equation}
For a partition $\mathcal R$ put
$C_{\mathcal R}u=(q_B(u))_{B\in\mathcal R}$, with the Euclidean product norm.
Its output dimension is $h\#\mathcal R$; on the full field space its rows
are independent, so its rank equals that dimension. A coefficient-only
partition is fixed by the coefficient profile and deterministic geometric
conventions before the field $u$ is chosen.

\subsection{An arbitrary symmetric matrix theorem}
The constants used below are intentionally conservative:
\begin{align}\label{eq:constants}
 b&=2^D,\quad J=2D,\quad
 C_{\rm ext}=1+J\{2J(1+12J)+1\},\nonumber\\
 G_D&=2^{2D+10}D^2C_{\rm ext}^2,\qquad
 B_D=(2b+3)bG_D^P,\\
 \alpha_D&=\frac D2\left(\frac D{D-2}\right)^{D-1}
                      \left(\frac34\right)^D.\nonumber
\end{align}
Here $G_D$ is the scalar weighted-gap constant and $B_D$ is the resulting
tree-count constant. Let $\lambda_j(x)$, $1\le j\le r$, be the eigenvalues
of $A(x)$, counted with multiplicity. Let $\nu_0$ count all its nonpositive
local eigenvalues over all $x$. Define
\begin{equation}\label{eq:cap}
 F_c(A)=\nu_0+
 \sum_{x,j:\lambda_j(x)>0}
       \min\{1,\alpha_D c^P\lambda_j(x)^{-P}\}.
\end{equation}
These are local coefficient eigenvalues, not the eigenvalues of the global
operator associated with $Q_A$.

\begin{theorem}[Capped-budget rectangular observations]\label{thm:matrix}
For every profile $A$ in \eqref{eq:actual} there is a coefficient-only
partition $\mathcal R$ into integer rectangles such that
\begin{equation}\label{eq:matrix-count}
 h\#\mathcal R\le
 2Dh\left[1+B_D+
       \max\left\{\frac{B_D}{\alpha_D},D+1\right\}F_c(A)\right].
\end{equation}
There is a finite coefficient-dependent $K$ such that, simultaneously for all
$u:X\to\R^h$,
\begin{equation}\label{eq:matrix-coercive}
 V^{-1}Q_A(u)+K\|C_{\mathcal R}u\|^2\ge c\normn{u}^2.
\end{equation}
One may take $K=c+4DL$, with $L$ explicitly defined in
\eqref{eq:negative-depth}. Neither the partition nor $K$ depends on $u$.
\end{theorem}

The budget caps the cost of each sufficiently small positive channel at one
and assigns the same unit cost to each nonpositive channel. Its count thus
does not grow with negative depth. The penalty $K$ can grow with that depth;
\eqref{eq:matrix-coercive} makes no assertion that $K=c$ for signed forms.
The elementary singleton partition has $hV$ outputs and also admits a finite
penalty for every finite symmetric form. One may choose whichever of that
partition and Theorem~\ref{thm:matrix} has fewer outputs. No optimal dimension,
sharp coefficient, or efficient implementation is asserted.

\subsection{The cost of a common full-vector partition}
The dimension in Theorem~\ref{thm:matrix} concerns a common partition with
all $h$ component means on each spatial member. The following refinement
shows why this measurement class can require a product of $h$ and the budget.

\begin{theorem}[Necessary product order for common partitions]\label{thm:common-rank}
Fix $D\ge3$ and $c>0$. For every pair of integers $m,h\ge1$ there is an actual positive
diagonal matrix profile on $X_n$, $n=3m+3$, with exactly $2Dm$ weak
direction-labelled channels and
\[
 2Dm\le F_c(A)\le2Dm+1.
\]
If a common partition $\mathcal R$ outputs all $h$ means
\eqref{eq:rows} and satisfies \eqref{eq:matrix-coercive} for every general
field $u:X_n\to\R^h$ with any finite quadratic penalty coefficient, then
\begin{equation}\label{eq:product-lower}
 \rank C_{\mathcal R}=h\#\mathcal R\ge hm
                  \ge\frac{hF_c(A)}{2D+1}.
\end{equation}
Thus no bound $C_D(h+F_c(A))$ holds uniformly in $h$ in this class:
setting $h=m$ makes the ratio of necessary rank to $h+F_c(A)$ diverge.
Together with Theorem~\ref{thm:matrix}, this gives matching worst-case
$hF_c(A)$ order, up to dimensional constants, in joint $h,F_c(A)$ growth.
The lower family consists of general vector fields and is not claimed to
be coexact or divergence-free. It does not address component-dependent
partitions, channel-adapted means, or postprocessed measurement classes.
\end{theorem}
\subsection{The scalar endpoint and the precise lower-bound domain}
For a positive scalar profile let
\begin{equation}\label{eq:weak-def}
 \begin{aligned}
 \mu&=V^{-1}\sum_xa(x),&Y(x)&=\mu/a(x),\\
 W(a,\mu)&=\sup_{\tau>0}\tau^P\Prob(Y>\tau),&
 T(n,\mu)&=\max\{1,n^D/(1+\mu)^P\},
 \end{aligned}
\end{equation}
where $\Prob$ is normalized counting probability. The strict superlevel
convention in this definition is used throughout.

\begin{theorem}[Weak endpoint upper bound]\label{thm:weakupper}
Suppose $1\le a\le2\mu$ and $W(a,\mu)\le M_w$, with $\mu$ the actual
arithmetic mean. For every $h\ge1$ there is a coefficient-only rectangular
partition with
\begin{equation}\label{eq:weak-upper-count}
 h\#\mathcal R\le
 2Dh\left\{1+2^Pc^PB_D[1+P M_w\log\mu]\right\}T(n,\mu)
\end{equation}
and
\begin{equation}\label{eq:scalar-coercive}
 V^{-1}Q_a(u)+c\|C_{\mathcal R}u\|^2\ge c\normn{u}^2
 \quad(u:X\to\R^h).
\end{equation}
\end{theorem}

To state the lower bound without an unspecified Hodge domain, fix $g\ge1$.
A flat one-cochain is $U=(U_1,\ldots,U_D)$ with $U_j:X\to\R^g$;
we regard it as an $h=Dg$ component field in $Q_a$.
For antisymmetric two-cochains $B_{ij}=-B_{ji}:X\to\R^g$, define
\begin{equation}\label{eq:coexact}
 D_i^-f(x)=f(x)-f(x-e_i),\qquad
 (\delta B)_j=-\sum_iD_i^-B_{ij},\qquad
 \mathcal H_{\rm coex}(X)=\{\delta B:B\text{ antisymmetric}\}.
\end{equation}
Backward divergence $\sum_jD_j^-U_j$ vanishes on this space by commutation
and antisymmetry. The energy on it is precisely the scalar component-gradient
$Q_a(U)$, with all its original forward starts; no weighted curl identity
is used to define or replace that energy.

\begin{theorem}[Weak endpoint obstruction and worst-case order]\label{thm:weaklower}
Fix $D\ge3$, $c>0$ and $g\ge1$. For every integer $m\ge1$ there are $n_m$,
a profile $a_m$ with actual mean $\mu_m$, and an $m$-dimensional real
subspace $E_m\subset\mathcal H_{\rm coex}(X_{n_m})$ such that
\begin{align}\label{eq:lower-data}
 1\le a_m\le2\mu_m,\qquad W(a_m,\mu_m)&\le1,\nonumber\\
 Q_{a_m}(U)&\le(c/4)\sum_x\|U(x)\|^2\quad(U\in E_m),\\
 T(n_m,\mu_m)&\le C(D,c),\qquad
 \Expect(\mu_m/a_m)^P\ge\frac{m}{2^P16^D}.\nonumber
\end{align}
Suppose an observation $C$ from a neighborhood of zero in
$\mathcal H_{\rm coex}(X_{n_m})$ to a finite-dimensional real normed space
$F$ is differentiable at zero and $C(0)=0$. Suppose also that
$\Phi(y)\le K\|y\|^2$ near zero for some finite $K$, and that locally
\begin{equation}\label{eq:local-coercive}
 c\sum_x\|U(x)\|^2\le Q_{a_m}(U)+\Phi(C(U)).
\end{equation}
Then $DC(0)|_{E_m}$ is injective and $\dim F\ge m$.
For this fixed-parameter sequence $m$ is comparable to $1+\log\mu_m$ for
large $m$. Thus the logarithmic loss in Theorem~\ref{thm:weakupper} has
matching worst-case order. This is not a claim of optimal rank for every
individual profile.
\end{theorem}

The centered, differentiable and local quadratic assumptions are part of
the lower theorem. Its penalty need not have a uniform bound as $m$ varies.
A common division of energy and squared field norm by $V$ gives the same
rank conclusion. The full critical moment grows along the lower family,
so this family does not contradict the strong-moment scalar theorem.

\subsection{Relationship to earlier work and proof roadmap}
This manuscript is a proposed linked successor to the AIRR record
\texttt{ARR-2026-7EPM3V3DP796BBRR}, version~13.
It combines the complete arbitrary-positive-scalar construction from that
predecessor with the subsequent matrix and weak-endpoint arguments. It is a
new source artifact; prior judgments do not transfer to it.

The closest constructive predecessors are dyadic-ring approximation and
compressed-tree arguments in Brudnyi~\cite{brudnyi221}, Sections~1--2 and
Appendices~I--II, and the bivariate BV approximation setting of
Cohen--DeVore--Petrushev--Xu~\cite{cdpx223}. Their approximation spaces and
objectives differ from the specified finite coefficient-only observation
problem here. This comparison is not a proof that the present statements
are absent from, or inequivalent to, the classical theory. In particular,
ring decomposition, boundary-touching holes, BV/Sobolev estimates and tree
compression are classical principles. The appendix spells out their finite
integer implementation, including heavy atoms and the original edge weights.

Frank~\cite{frank}, Theorems~2.1 and~3.1, supplies the relevant classical
Cwikel/CLR and operator-valued spectral context. A weak multiplier hypothesis
in such a theorem must be distinguished from replacing the strong coefficient
moment in the present problem by a weak spatial distribution hypothesis.
Our matrix proof uses elementary local spectral calculus and a scalar minimum
eigenvalue bound; it does not invoke a global CLR estimate. The particular
positive normalization $\alpha_D$ in \eqref{eq:constants} fixes the cap
threshold and numerical constants; the argument works with any fixed positive
cap normalization after corresponding changes. Spectral replacement,
scalar domination, rectangle refinement, layer cake and disjoint low-mode
rank arguments are classical. No new general spectral principle or
historical priority is claimed.

Section~\ref{sec:matrixproof} proves Theorem~\ref{thm:matrix} using the
fully proved scalar theorem in Appendix~\ref{sec:balanced223}.
Section~\ref{sec:common-rank} proves the common-partition product lower bound.
Section~\ref{sec:endpoint} gives the layer-cake upper bound and the complete
compact-bump coexact lower construction. The appendix proves the scalar
BV extension, actual weighted ring gap, heavy-chain count and rectangular
mean inequality without an assumed geometric or local-gap certificate.

All assertions concern finite unit-spacing component-gradient forms.
No interacting or covariant coefficient producer, variable-weight curl
replacement, continuum theorem or Yang--Mills solution follows. The programme
context PR29=$20/41$, TermSource $0$, is unchanged.

\section{Positive replacement and exact rectangular refinement}\label{sec:matrixproof}
We prove Theorem~\ref{thm:matrix}. Every step uses the actual matrix profile
on the forward stack in \eqref{eq:actual}. In particular, positivity is
used to discard wrap terms only after replacing the signed profile.

\subsection{The spectral split and scalar inverse moment}
Set
\[
 \sigma=\alpha_D^{1/P}c,
 \qquad S_x=1_{(-\infty,\sigma]}(A(x))\R^r,
 \qquad N=\sum_x\dim S_x,\qquad I=\{x:S_x\ne\{0\}\}.
\]
Here $1_{(-\infty,\sigma]}(A(x))$ is the spectral projector. A degenerate
eigenspace at the cutoff is selected in full, so no choice of basis within
it enters the construction. Every nonpositive channel is selected. A positive
selected channel contributes exactly one to the capped budget, and an
unselected channel contributes its uncapped inverse power. Consequently
\begin{equation}\label{eq:cap-exact}
 F_c(A)=N+\alpha_Dc^P
                \sum_{x,j:\lambda_j(x)>\sigma}\lambda_j(x)^{-P},
 \qquad 0\le N\le F_c(A),\quad |I|\le N.
\end{equation}
Equality $\lambda=\sigma$ is included in $N$.

Put $T=\sigma+c(1+N)>\sigma$. Replace every selected eigenvalue by $T$,
leaving all unselected eigenvalues and every local spectral subspace
unchanged, and denote the resulting positive definite matrix by $A_T(x)$.
For $N=0$ this changes nothing. With
$a_T(x)=\lambda_{\min}(A_T(x))>0$, pointwise matrix positivity gives
\begin{equation}\label{eq:scalar-domination}
 Q_{A_T}(u)\ge Q_{a_T}(u).
\end{equation}
This controls the complete stacked vector, including all directional and
internal mixing. It requires no commutation between distinct matrices.
Moreover, the largest eigenvalue of $A_T(x)^{-P}$ is at most its trace, hence
\begin{align}\label{eq:replacement-moment}
 c^P\sum_xa_T(x)^{-P}
 &\le c^P\sum_x\tr(A_T(x)^{-P})\nonumber\\
 &=c^P\sum_{\lambda_j(x)>\sigma}\lambda_j(x)^{-P}
                  +N(c/T)^P\\
 &\le \frac{F_c(A)-N}{\alpha_D}+1.\nonumber
\end{align}
The last inequality uses $T\ge c(1+N)$ and
$N/(1+N)^P\le1$ for $P\ge1$. When $N=0$ the extra one is harmless.

\subsection{The scalar partition is derived, not assumed}
Theorem~\ref{thm:balanced223} in Appendix~\ref{sec:balanced223} applies to
every finite positive $a_T$, without an arithmetic-mean, lower-envelope or
upper-envelope condition. Its partition $\mathcal R_0$, combined with
\eqref{eq:scalar-domination}, satisfies
\begin{equation}\label{eq:positive-observation}
 c\normn{u}^2\le V^{-1}Q_{A_T}(u)+c\|C_{\mathcal R_0}u\|^2,
\end{equation}
by summing the same scalar inequality over all $h$ real components.
If $R_0=\#\mathcal R_0$, its exact count and
\eqref{eq:replacement-moment} yield
\begin{align}\label{eq:initial-count}
 R_0&\le2D\left[1+(2b+3)b(cG_D)^P\sum_xa_T(x)^{-P}\right]\nonumber\\
 &\le2D\left[1+B_D\left(1+\frac{F_c(A)-N}{\alpha_D}\right)\right].
\end{align}
There is no auxiliary local-gap or extension hypothesis in this application:
the scalar construction and its gap are proved in the appendix.

\subsection{Isolating every affected endpoint}
Let
\begin{equation}\label{eq:endpoints}
 Z=I\cup\bigcup_{i=1}^D(I+e_i),\qquad |Z|\le(D+1)|I|\le(D+1)N,
\end{equation}
where the translations are periodic. We refine $\mathcal R_0$ so that each
$z\in Z$ is a singleton member. Order $Z$ lexicographically in the fixed
chart. If $z$ is not already a singleton, let $R$ be its current rectangle.
Partition $R\setminus\{z\}$ according to the first coordinate in which a
point differs from $z$, and according to whether it is smaller or larger in
that coordinate. Earlier coordinates are fixed to those of $z$, later
coordinates range over their original intervals, and the differing coordinate
ranges below or above $z$. Each nonempty piece is an integer rectangle.
There are at most $2D$ such pieces. Adding $\{z\}$ replaces one rectangle
by at most $2D+1$, increasing the count by at most $2D$.
Earlier isolated points lie in disjoint singleton members and remain intact.
Thus the final partition $\mathcal R$ satisfies
\begin{equation}\label{eq:refined-count}
 R:=\#\mathcal R\le R_0+2D|Z|\le R_0+2D(D+1)N.
\end{equation}
A wrapping forward endpoint is simply its chart representative in $Z$;
this causes no alteration of the rectangle argument.

For each original parent $B$ and its final pieces $B_k$,
\[
 q_B(u)=\sum_k\sqrt{|B_k|/|B|}\,q_{B_k}(u),\qquad
 \sum_k|B_k|/|B|=1.
\]
Hilbert-space Cauchy with these exact weights proves
\begin{equation}\label{eq:refine-monotone}
 \|C_{\mathcal R_0}u\|^2\le\|C_{\mathcal R}u\|^2,
 \qquad V^{-1}\sum_{z\in Z}\|u(z)\|^2
                      \le\|C_{\mathcal R}u\|^2.
\end{equation}
The second assertion holds because each $z\in Z$ contributes the singleton
row $V^{-1/2}u(z)$. The remaining rows are still only component means on
rectangles; no gradient or nonrectangular observation has been introduced.

\subsection{Charging the complete signed correction}
Define the finite number
\begin{equation}\label{eq:negative-depth}
 L=\begin{cases}
 0,&N=0,\\
 \displaystyle\max_{x,j:\lambda_j(x)\le\sigma}(T-\lambda_j(x)),&N>0.
 \end{cases}
\end{equation}
The spectral construction gives
$0\le A_T(x)-A(x)\le L\Id$, with zero difference for $x\notin I$.
This is an operator inequality on all $r$ coordinates; in particular it
bounds all rotated off-diagonal contributions together. It does not assign
signs to individual cross terms. Therefore
\begin{align}\label{eq:correction}
 0\le Q_{A_T}(u)-Q_A(u)
 &\le L\sum_{x\in I}\|du(x)\|^2\nonumber\\
 &\le2L\sum_{x\in I}
      \left[D\|u(x)\|^2+\sum_i\|u(x+e_i)\|^2\right]\\
 &\le4DL\sum_{z\in Z}\|u(z)\|^2.\nonumber
\end{align}
Indeed each site has its own bad-start multiplicity at most $D$ and at most
$D$ incoming direction-labelled starts. All their endpoints belong to $Z$.
This counts parallel edges at period two separately, as required by the
actual form. At $n=1$ every difference vanishes and the estimate remains
valid. No periodic term has been dropped from the signed correction.

Insert \eqref{eq:correction} into \eqref{eq:positive-observation}, divide
the point sum by $V$, and use \eqref{eq:refine-monotone}. This proves
\eqref{eq:matrix-coercive} with $K=c+4DL$. The count is
\begin{align*}
 hR&\le2Dh\left[1+B_D+\frac{B_D}{\alpha_D}(F_c(A)-N)
                                      +(D+1)N\right]\\
 &\le2Dh\left[1+B_D+
                 \max\{B_D/\alpha_D,D+1\}F_c(A)\right],
\end{align*}
using the nonnegativity of both $N$ and $F_c(A)-N$. This completes the proof
of Theorem~\ref{thm:matrix}.

\begin{remark}[What the cap does and does not control]\label{rem:cap}
The only use of the budget after the local spectral split is the exact
identity \eqref{eq:cap-exact}. In particular its control of observation
number is compatible with arbitrarily negative selected eigenvalues, since
the penalty absorbs the actual finite correction. When $N=0$, the correction
vanishes and $K=c$. For a positive scalar profile the sharper direct scalar
construction also has $K=c$ without selecting or replacing any channels.
For completeness, if $A(x)\ge-M\Id$ for all $x$, then
$Q_A(u)\ge-4DM\sum_x\|u(x)\|^2$, because
$\sum_x\|du(x)\|^2\le4D\sum_x\|u(x)\|^2$.
For the singleton partition $\|C_{\mathcal R}u\|^2=\normn{u}^2$,
so $K=c+4DM$ proves its asserted finite-penalty alternative.
The matrix construction establishes an existence bound with explicit
constants, rather than a sharp spectral rank formula.
\end{remark}

\section{Necessary product order for full-vector common partitions}\label{sec:common-rank}
We prove Theorem~\ref{thm:common-rank} by an explicit positive diagonal
family. The proof concerns all vector fields, as in Theorem~\ref{thm:matrix},
rather than the smaller coexact domain of the endpoint lower theorem.
It uses only local fields and finite-dimensional rank-nullity.

\subsection{The coefficient family and every local channel}
Let $n=3m+3$ and choose centers
\[
 x_j=(3j+1,1,\ldots,1),\qquad 0\le j<m.
\]
They lie away from the periodic wrap boundary. With
$\sigma=\alpha_D^{1/P}c$, define
\[
 a_0=\min\{c/(8D),\sigma/2\}>0,\qquad
 M=r n^D,\qquad H=\sigma M>\sigma.
\]
In the canonical direction/component basis $(i,k)$ of $\R^{Dh}$ let
$A(x)$ be diagonal, with entries
\begin{equation}\label{eq:rank-profile}
 \lambda_{i,k}(x)=
 \begin{cases}
 a_0,&k=1\text{ and }x\in\{x_j,x_j-e_i\}\text{ for some }j,\\
 H,&\text{otherwise}.
 \end{cases}
\end{equation}
At a center all $D$ direction channels for component one are weak; at an
incoming start only its specified direction/component channel is weak.
All eigenvalues are positive. For each fixed direction the starts $x_j$
and $x_j-e_i$ are distinct. Across centers their first coordinates lie in
disjoint sets $\{3j,3j+1\}$. Channels with different direction labels remain
distinct even when they share a start. There are therefore exactly $2Dm$
weak channels.

Since $a_0\le\sigma/2$, each weak channel has capped cost one in
\eqref{eq:cap}. Every other channel has capped cost
$\alpha_Dc^PH^{-P}=M^{-P}$. Thus, including all strong channels,
\begin{equation}\label{eq:rank-budget}
 F_c(A)=2Dm+(M-2Dm)M^{-P},\qquad
 2Dm\le F_c(A)\le2Dm+1.
\end{equation}
Indeed $0\le M-2Dm\le M$ and $P>1$, so the tail is at most
$M^{1-P}\le1$. Increasing $h$ increases $H$ and keeps this strong tail
bounded; it introduces no uncontrolled $h$-sized budget contribution.
The choice of $H$ is permitted by the arbitrary finite profile hypothesis.

\subsection{The actual low space and the restriction rank}
Let $U_j(x)=1_{\{x_j\}}(x)e_1$, where $e_1\in\R^h$ is the first
canonical unit vector. These fields are orthonormal in counting units.
Their only nonzero direction-labelled forward differences are $-e_1$ at
$x_j$ and $+e_1$ at $x_j-e_i$, for each direction $i$. Both original
starts are assigned coefficient $a_0$ by \eqref{eq:rank-profile}; there
are no further wrap contributions. Consequently
\[
 Q_A(U_j)=2Da_0\le c/4.
\]
No nearest-neighbor edge joins distinct centers, and their gradient-start
supports are disjoint. Hence the quadratic cross terms vanish and
\begin{equation}\label{eq:rank-lowspace}
 L_m=\operatorname{span}_{\R}\{U_0,\ldots,U_{m-1}\},\quad
 \dim L_m=m,\quad
 Q_A(U)\le(c/4)\sum_x\|U(x)\|^2\quad(U\in L_m).
\end{equation}
These are general vector functions. No divergence or coexact property is
asserted for them.

For any common partition $\mathcal R$ into $R$ nonempty rectangles, its
actual rows on this space are
\[
 q_B\left(\sum_jt_jU_j\right)
 =V^{-1/2}|B|^{-1/2}
       \left(\sum_{j:x_j\in B}t_j\right)e_1.
\]
Each rectangle supplies at most one scalar functional on $L_m$, so
$\rank(C_{\mathcal R}|_{L_m})\le R$. If $R<m$, rank-nullity gives a
nonzero $U\in L_m$ with $C_{\mathcal R}U=0$. Coercivity with any finite
quadratic penalty, after multiplication by $V$, would imply
\[
 c\sum_x\|U(x)\|^2\le Q_A(U)
                    \le(c/4)\sum_x\|U(x)\|^2,
\]
a contradiction. Thus every successful common partition has $R\ge m$.
The same argument works for any penalty on these linear observations with
$\Phi(0)=0$, without a differentiability assumption on that penalty.

The full field-space rank is exactly $hR$, not merely a padded output
dimension. To see this, for each $B$ choose a site $z\in B$. A field
supported at $z$ with value $\sqrt{V|B|}\,e_k$ has observation $e_k$ in
row $B$ and zero in every other row. These fields show surjectivity onto
all $hR$ row coordinates. Therefore necessary genuine rank is at least
$hm$. Since $F_c(A)\le2Dm+1\le(2D+1)m$, this proves
\eqref{eq:product-lower}.

\subsection{Joint growth and the measurement-class restriction}
Set $h=m$. The necessary rank is at least $m^2$, while
$h+F_c(A)\le(2D+1)m+1$. Their ratio tends to infinity for fixed $D,c$.
Thus a dimension-only $C_D(h+F_c(A))$ bound fails for these common
full-vector partitions, even with arbitrary finite penalty. The product
lower bound and Theorem~\ref{thm:matrix} have matching $hF_c(A)$ order
along this family, since $F_c(A)\ge2Dm$ makes $1+F_c(A)$ comparable to
$F_c(A)$. Sharp coefficients, additive constants and per-profile optimality
are not determined.

The extra factor $h$ measures the cost of retaining every component mean
on every common partition member; it does not count $h$ independent tests
of the one-component low space. The argument does not apply to distinct
componentwise partitions, projected or channel-adapted rows, or compressed
postprocessed outputs. It requires no external theorem about the rank of
those broader observation classes. These elementary rank and local-field
arguments carry no claim of historical priority, a new general principle,
formal verification, or physical transfer.

\section{The weak critical endpoint}\label{sec:endpoint}
The positive and negative results use the same scalar original-start energy.
The positive estimate holds on every component field. The negative estimate
is already forced on the precisely defined coexact subspace
\eqref{eq:coexact}.

\subsection{Finite layer cake and the upper estimate}
Since $a\ge1$, its actual mean obeys $\mu\ge1$ and $0<Y\le\mu$.
For each value $y$ in this finite distribution,
\[
 y^P=P\int_0^\mu\tau^{P-1}1_{\{\tau<y\}}\,d\tau.
\]
Summing over $X$ proves the exact identity
\begin{equation}\label{eq:layercake}
 \Expect Y^P=P\int_0^\mu\tau^{P-1}\Prob(Y>\tau)\,d\tau
                   \le1+P M_w\log\mu.
\end{equation}
The contribution on $(0,1)$ is at most one, and on $[1,\mu]$ we use
$\Prob(Y>\tau)\le M_w\tau^{-P}$. At $\mu=1$ the second integral is zero;
in fact the actual mean and $a\ge1$ then force $a\equiv1$.
Atoms at strict cutoffs change no ordinary integral. The envelope
$a\le2\mu$ is retained in the theorem but is not needed for this step.

Corollary~\ref{cor:balanced223-moment} in the appendix, with the actual
moment $M=\Expect Y^P$, gives
\[
 h\#\mathcal R\le2Dh[1+2^Pc^PB_DM]T(n,\mu),
\]
and the normalized inequality \eqref{eq:scalar-coercive} with penalty $c$.
Substitution of \eqref{eq:layercake} proves
Theorem~\ref{thm:weakupper}. Its geometry and gap are supplied by the
complete scalar construction, rather than by a separate assumed certificate.

\subsection{A compact product bump on the actual coexact domain}
We now prove Theorem~\ref{thm:weaklower}. Fix a unit vector $e\in\R^g$.
Define a globally $C^1$ function
\[
 \beta(s)=\begin{cases}s^2(1-s)^2,&0\le s\le1,\\0,&s\notin[0,1].\end{cases}
\]
We use $\beta$ for the bump to distinguish it from the critical exponent
$P$. Both $\beta$ and $\beta'$ vanish at the endpoints. On $[0,1]$,
\[
 \beta'(s)=2s(1-s)(1-2s),\qquad
 \beta''(s)=2-12s+12s^2.
\]
Thus $0\le\beta\le1$, $|\beta'|\le2$, and $\beta'$ is globally
$2$-Lipschitz: its interior second derivative has absolute value at most two,
and the zero extension has no derivative jump.

For an integer $L\ge16$ divisible by $16$, put
\[
 \psi_L(x)=\prod_{i=1}^D\beta(x_i/L),\qquad x\in\Z^D,
\]
and define the one-cochain
\begin{equation}\label{eq:bumpfield}
 U_1=D_2^-\psi_L\,e,\qquad U_2=-D_1^-\psi_L\,e,
 \qquad U_j=0\quad(j\ge3).
\end{equation}
Taking $B_{12}=\psi_Le$, $B_{21}=-\psi_Le$ and all other $B_{ij}=0$
in \eqref{eq:coexact} gives exactly \eqref{eq:bumpfield}, including its
signs. Therefore $U=\delta B$ and
\[
 \sum_jD_j^-U_j=D_1^-D_2^-\psi_Le-D_2^-D_1^-\psi_Le=0.
\]
Every component has zero mean by telescoping. The same identities hold
exactly on any torus large enough to contain this support away from wraps.

The one-dimensional lattice sequence $f(t)=\beta(t/L)$ is supported at
$t=1,\ldots,L-1$, and $D^-f$ is supported at $t=1,\ldots,L$.
It follows from the product formula that $U$ is supported in
$\{0,\ldots,L\}^D$ and vanishes on every lower coordinate face $x_i=0$.
If $D_i^+U_j(x)\ne0$, its original start belongs to this same box.
Indeed the only potentially external incoming start is $x_i=-1$; its next
endpoint is on the zero lower face, so both field values vanish. Upper
outgoing starts $x_i=L$ can contribute and are included. In all other
coordinates an external coordinate makes both terms zero. This verifies the
boundary-start condition needed below with no interior-only energy convention.

\begin{lemma}[Finite bump constants]\label{lem:bump}
Let $\E_1(U)=\sum_{x,i,j}\|D_i^+U_j(x)\|^2$ on $\Z^D$. Then
\begin{align}\label{eq:bumpbounds}
 \E_1(U)&\le U_D L^{D-4},& U_D&=32D\,2^D,\nonumber\\
 \sum_x\|U(x)\|^2&\ge\eta_D L^{D-2},&
 \eta_D&=\frac9{2^{D+12}256^{2(D-1)}}.
\end{align}
In particular, with $K_D=U_D/\eta_D$,
\begin{equation}\label{eq:bumprayleigh}
 \E_1(U)\le (K_D/L^2)\sum_x\|U(x)\|^2.
\end{equation}
\end{lemma}
\begin{proof}
The Lipschitz bound on $\beta'$ bounds each same-coordinate second finite
difference by $2/L^2$. For different coordinates the product of two first
differences is bounded by $4/L^2$. All unused factors have magnitude at most
one. At most $2D$ component/direction pairs are nonzero; each has at most
$(L+1)^D\le2^DL^D$ original starts by the support analysis. The common
squared bound $16/L^4$ gives the first line of \eqref{eq:bumpbounds},
including the upper outgoing starts.

For the norm, take $x_2=L/8+1,\ldots,L/4$, exactly $L/8$ sites,
and $x_i\in[L/4,3L/4]$ for every $i\ne2$, at least $L/2$ sites each.
On each backward step in coordinate two, the mean-value point belongs to
$[1/8,1/4]$, where $\beta'\ge3/32$. On the other coordinates
$\beta\ge1/256$ on $[1/4,3/4]$. Hence
\[
 \|U_1(x)\|\ge\frac3{32L}256^{-(D-1)}
\]
on all these sites. Squaring and multiplying by
$(L/8)(L/2)^{D-1}$ gives the stated $\eta_D$. It is strictly positive,
so division yields \eqref{eq:bumprayleigh}.
\end{proof}

\subsection{Geometric scales, exact mean and the weak distribution}
Choose an integer $k\ge0$ such that $\varepsilon=4^{-k}$ satisfies
$\varepsilon K_D\le c/4$. This choice depends on $D,c$ and is fixed
while $m$ varies. Put
\begin{equation}\label{eq:scales}
 L_0=2^{\max\{4,k\}},\quad L_j=2^jL_0\ (0\le j<m),\quad
 L_*=L_{m-1},\quad n=16L_*,\quad H=\varepsilon L_*^2.
\end{equation}
Every $L_j$ is divisible by $16$, and $\varepsilon L_0^2\ge1$.
Place boxes
\[
 B_j=q_j+\{0,\ldots,L_j\}^D,\quad
 q_{0,1}=1,\quad q_{j+1,1}=q_{j,1}+L_j+3,\quad
 q_{j,i}=1\ (i\ge2).
\]
Their first-coordinate intervals have two unused sites between successive
boxes. The rightmost endpoint is $\sum_jL_j+3m-2$.
Since $\sum_jL_j<2L_*$ and $m\le2^{m-1}\le L_*/16$,
this is less than $3L_*+1<n$. Other coordinates also lie strictly inside
the chart. Thus no support or gradient-start support collides periodically.
Translate \eqref{eq:bumpfield} into each $B_j$, obtaining $U^{(j)}$.
Their field supports and their complete original gradient-start supports
are disjoint.

Define the actual scalar coefficient by
\begin{equation}\label{eq:weakprofile}
 a(x)=\begin{cases}\varepsilon L_j^2,&x\in B_j,\\
 H,&x\notin\bigcup_jB_j.
 \end{cases}
\end{equation}
The boxes are disjoint and the top box has the same coefficient as the
background. In particular $1\le a\le H$. Define its mean by its actual sum:
\begin{equation}\label{eq:exactmean}
 \mu=V^{-1}\sum_xa(x)
    =H-V^{-1}\sum_j|B_j|(H-\varepsilon L_j^2).
\end{equation}
This is not a freely adjusted normalization. With
\[
 C_{\rm vol}=\frac{2^D}{1-2^{-D}},
\]
the geometric series and $(L_j+1)^D\le2^DL_j^D$ imply
\[
 \sum_j|B_j|\le C_{\rm vol}L_*^D,\qquad
 V^{-1}\sum_j|B_j|\le C_{\rm vol}/16^D<1/2.
\]
It follows that
\begin{equation}\label{eq:meanbounds}
 H/2\le\mu\le H,\qquad \mu\ge1,\qquad 1\le a\le H\le2\mu.
\end{equation}
The separate inequality $\mu\ge1$ follows directly from $a\ge1$,
including when $H/2<1$.

We check $W(a,\mu)\le1$ with strict cutoffs. If
$0<\tau\le\mu/H$, then $\tau\le1$ and
$\tau^P\Prob(Y>\tau)\le1$. If $\tau>\mu/H$, the background and top
box are absent. The boxes contributing to $Y>\tau$ satisfy
$\varepsilon L_j^2<\mu/\tau$. If there is any such box, let $j$ be the
largest qualifying index. They form an initial segment and have total
volume at most $C_{\rm vol}L_j^D$. Consequently
\begin{align}\label{eq:weakcheck}
 \tau^P\Prob(Y>\tau)
 &\le C_{\rm vol}(\mu/\varepsilon)^P/V\nonumber\\
 &\le C_{\rm vol}L_*^D/(16L_*)^D<1.
\end{align}
If no box qualifies the probability is zero. Equality at a coefficient atom
excludes that box under the strict criterion and cannot increase the bound.
This proves the weak distribution assertion for every $m$, including $m=1$.

\subsection{The low space and every finite local quadratic penalty}
Every nonzero gradient of $U^{(j)}$ has its original start in $B_j$,
where \eqref{eq:weakprofile} has coefficient $\varepsilon L_j^2$.
All wrap terms vanish for these nonwrapping supports. Thus the complete
energy, not merely its interior contribution, is
\begin{equation}\label{eq:weightedbump}
 Q_a(U^{(j)})=\varepsilon L_j^2\E_1(U^{(j)})
 \le\varepsilon K_D\sum_x\|U^{(j)}(x)\|^2
 \le(c/4)\sum_x\|U^{(j)}(x)\|^2.
\end{equation}
The norm bound in Lemma~\ref{lem:bump} makes every $U^{(j)}$ nonzero.
Disjoint field supports give norm orthogonality; disjoint gradient-start
supports give vanishing quadratic-energy cross terms. The same inequality
therefore holds on
\[
 E_m=\operatorname{span}_{\R}\{U^{(0)},\ldots,U^{(m-1)}\}
       \subset\mathcal H_{\rm coex}(X),\qquad \dim E_m=m.
\]
This proves an actual low subspace, with no local eigenfunction or gap
assumption. From \eqref{eq:scales} and \eqref{eq:meanbounds},
\begin{equation}\label{eq:targetbounded}
 T(n,\mu)\le\left(16\sqrt{2/\varepsilon}\right)^D.
\end{equation}
The right side is independent of $m$.

For clarity the finite-dimensional obstruction is proved directly. If
$DC(0)|_{E_m}$ has a nonzero kernel vector $U$, differentiability and
$C(0)=0$ imply $C(tU)=o(t)$. For sufficiently small $t$, the assumed
upper bound on the penalty gives
\[
 \Phi(C(tU))\le K\|C(tU)\|^2=o(t^2).
\]
Apply \eqref{eq:local-coercive} to $tU$, use the low-space inequality,
and divide by $t^2$. Letting $t\to0$ would imply
\[
 c\sum_x\|U(x)\|^2\le(c/4)\sum_x\|U(x)\|^2,
\]
a contradiction. Therefore the derivative is injective on $E_m$ and
$\dim F\ge\rank(DC(0)|_{E_m})=m$. No uniformity of $K$ between profiles
was used. Positivity or differentiability of $\Phi$ is unnecessary; only
its stated finite quadratic upper bound is used. Common normalization of
energy and field norm changes none of this finite-dimensional argument.

Finally, disjoint boxes, $|B_j|\ge L_j^D$ and $\mu\ge H/2$ give
\begin{equation}\label{eq:fullmomentgrowth}
 \Expect(\mu/a)^P
 \ge V^{-1}\sum_{j=0}^{m-1}L_j^D
                   \left(\frac\mu{\varepsilon L_j^2}\right)^P
 =\frac{m(\mu/\varepsilon)^P}{V}
 \ge\frac m{2^P16^D}.
\end{equation}
This establishes every part of \eqref{eq:lower-data} and explains why the
strong-moment theorem is not contradicted.

\subsection{The matching order and its quantifiers}
For fixed $D,c$, $\varepsilon$ and $L_0$ remain fixed. Equations
\eqref{eq:scales} and \eqref{eq:meanbounds} imply
\begin{equation}\label{eq:logmean}
 \log\mu=2(m-1)\log2+\log(\varepsilon L_0^2)+\theta_m,
 \qquad -\log2\le\theta_m\le0.
\end{equation}
Thus $m/(1+\log\mu)\to1/(2\log2)$.
The lower bound just proved forces output dimension at least a positive
fixed-parameter multiple of $1+\log\mu$ for large $m$.
Theorem~\ref{thm:weakupper}, with $h=Dg$, $M_w=1$ and
\eqref{eq:targetbounded}, gives an upper bound of the same order on these
profiles. In particular, no bound proportional to
$T(n,\mu)\,o(\log\mu)$ can hold uniformly over the weak-bound class for
the centered differentiable/local-quadratic observations under discussion.
This proves the final assertion of Theorem~\ref{thm:weaklower}.

The conclusion is a worst-case comparison across profiles. It does not
state that each profile needs a logarithm or identify its minimal rank.
The lower construction also does not cover a nonzero observation baseline
with an unrelated positive penalty at that baseline, nondifferentiable
encodings, or penalties outside the stated local quadratic class. Those
assumptions matter to the argument rather than merely to its presentation.

\section{Model assessments and review context}\label{sec:review}
Lluis Eriksson is the named author. This manuscript was drafted with
assistance from Codex, exposed as an OpenAI system based on GPT-6. The exact
served backend and reasoning-effort setting are unknown. This drafting
assistance is distinct from the involved AI source-audit roles described below.

At preparation, this exact manuscript has not yet been numerically assessed.
Its predecessor's judgments do not transfer to these source bytes or to a
future compiled PDF. The preparation uses prior source audits of the scalar,
weak-endpoint, capped-matrix and common-partition rank arguments. Those audits involved AI roles in
the same agent system, with parent framing and shared-context limits; the
true backend identities are unknown. They are not represented as blind,
independent or human reviews. The author of this manuscript is not its
source auditor or numerical reviewer.

The complete original source-audit reports, their adverse findings and scope
limits remain external companions, identified by raw-byte hashes in the
source-input manifest. Their qualified acceptances concern the named
arguments and do not certify priority, practical complexity, physical
transfer, optimal constants, all-size execution, or this manuscript's
presentation. No numerical grade is inferred from them. This manuscript
requires its own review of the exact source and any exact PDF produced from
it.

The selected diagnostics recorded for unit~224 used owned Colab Linux runs
in normal and optimized modes: six direct tensor-sum cases, one actual
$D=3$, $L=16$ coexact case, fifteen profiles and three rejected real negative
controls. Unit~225 recorded seven finite matrix families, 131 actual fields
and four rejected real controls, including endpoint, normalization and
full off-diagonal correction mutations. The normal and optimized outputs
were recorded as identical with zero actual exits. These are selected
diagnostics, not an all-size formal proof or a two-fresh-clone replay of
the construction. Their runtimes were released and read back inactive;
billed compute units are unknown. The mathematical proofs above and below
do not depend on these diagnostic counts.

\subsection*{Source identity and external evidence}
The scalar appendix transcribes the complete mathematical portion of
\texttt{balanced-rings223.tex} at programme commit
\begin{center}\small\texttt{ddd2cf605c1715a03d2e87991e920c0ec760e527}.\end{center}
Its raw Git-blob SHA-256, without checkout conversion or normalization, is
\begin{center}\small\ttfamily
6877841cde546f8920c39b98a5dfef8a\\
f385b1a2d5a4842f95a9dd36985b3dc2.
\end{center}
Changes to that mathematical portion are notation and context adaptations:
arbitrary $h$ real components are obtained by summing the same scalar
inequality; normalized and counting rows are distinguished; cross-references
are internal to this manuscript. The original extended operational and
assessment history is retained externally rather than interleaved with the
proof.

The full subsequent analytical sources are the unit~224 draft and its
logarithmic addendum, and the unit~225 capped-matrix draft. Their raw-byte
SHA-256 digests, respectively, are
\begin{center}\small\ttfamily
7a94460f2cbbaf6ac291d64981a7ae6e7\\
9590067f315923b615f842d6a0b211c\par\medskip
0720cedc459361539c9cedf95b7b1868\\
31a6ab9414ea27c859918203a39fb059\par\medskip
a9c4bcbc3b8764b2e2cf84b7906ae5fe9\\
ed3fa96c5f5ae22c68eb665a0f7d448.
\end{center}
The related unit~227 draft and complete source-only audit are also included
by identity in the external manifest. Its audit returned a qualified acceptance
for the stated common-full-vector measurement class. No unit~227 diagnostics
are claimed in this preparation; its lower proof uses no external spectral
$h+F$ theorem.

The accompanying \texttt{SOURCE-INPUTS.json} names the exact source files,
complete audit reports, context disclosures and unit result records with
their raw-file hashes. All prior reports remain unchanged. This evidence
identifies the inputs; it does not turn their recorded finite checks into
an executable proof for every size.

\appendix
\section{Coefficient-selected rectangular means for arbitrary scalar profiles}
\label{sec:balanced223}
This appendix gives the complete scalar construction used in the main proofs.
It uses the original positive-start scalar coefficient and
retains the edges crossing rectangular observation pieces inside a ring.
It does not identify a variable-weight curl/divergence form with that gradient
form. The balanced-ring, BV and tree mechanisms are classical;
historical originality of this finite specialization is unestablished.

\subsection{Statement and actual observations}
Let integers $D\ge3$, $n\ge1$, and $c>0$ be fixed, and put
$X=\{0,\ldots,n-1\}^D$. For any finite scalar profile $a(x)>0$ define
\[
 P=D/2,\quad b=2^D,\quad J=2D,\qquad
 C_{\rm ext}=1+J\{2J(1+12J)+1\},\qquad
 G_D=2^{2D+10}D^2C_{\rm ext}^{\,2}.
\]
Write $S(A)=\sum_{x\in A}a(x)^{-P}$ and
$t=(cG_D)^{-P}/b$. Every internal edge below keeps its original positive
start, and all sums and norms initially use counting measure. Here $\E_a(f;\Omega)$
denotes the sum over every chart-internal edge with both endpoints in $\Omega$,
weighted by its original positive-start value $a(x)$; $\E_1$ has unit weights.
Thus the counting row $\widehat q_B$ used at the end of this appendix equals
$V^{1/2}q_B$ in the normalized notation of \eqref{eq:rows}.
\begin{theorem}\label{thm:balanced223}
There is a partition $\mathcal R$ of $X$ into nonempty axis-aligned integer
rectangles, determined by $a$ and fixed geometric conventions alone, such that
\[
 \#\mathcal R\le2D\{1+(2b+3)S(X)/t\},
\]
and, simultaneously for every real field $f$,
\begin{equation}\label{eq:balanced223-final}
 \sum_{i,x:\ x,x+e_i\in X}a(x)|f(x+e_i)-f(x)|^2
 +c\sum_{B\in\mathcal R}\frac1{|B|}\Bigl(\sum_{x\in B}f(x)\Bigr)^2
 \ge c\sum_{x\in X}|f(x)|^2.
\end{equation}
The original periodic energy also satisfies this inequality, by adding
its nonnegative wrap terms. For every integer $h\ge1$ real components the same
partition has output dimension $h\#\mathcal R$ and penalty $\Phi(y)=c\|y\|^2$.
\end{theorem}
\begin{corollary}\label{cor:balanced223-moment}
For any reference value $\mu\ge1$, if
$n^{-D}\sum_x(\mu/a(x))^{D/2}\le M$, the actual normalized rectangular
means have output dimension at most $C_{\rm rect}T(n,\mu)$, with $K=c$ and
\[
 C_{\rm rect}=2Dh\{1+(2b+3)b(cG_D)^P2^PM\}.
\]
The exact arithmetic mean and envelope $1\le a\le2\mu$ are unnecessary
for this sufficiency. The inverse-moment hypothesis is still assumed;
it is not deduced from an arithmetic mean.
\end{corollary}
\begin{proof}
$S(X)\le Mn^D\mu^{-P}\le2^PM\max\{1,n^D/(1+\mu)^P\}$, since
$\mu\ge1$. Insert $1/t=b(cG_D)^P$ into the theorem. Apply its scalar
inequality to each of the $h$ real components. Dividing energy and squared
norms by $V=n^D$ gives observation rows
$V^{-1/2}|B|^{-1/2}\sum_B f=\sqrt{|B|/V}\,\operatorname{mean}_B f$.
The all-field inequality restricts to divergence-zero fields, and linear
means with the displayed quadratic penalty supply the original local data.
\end{proof}

\subsection{Balanced integer nodes and actual collars}
Split every coordinate interval with at least two sites into its floor/ceiling
halves; a singleton remains unchanged. Use coordinate order $1,\ldots,D$,
lower before upper, and lexicographic order on children. Each node has at
most $b$ children, and all side lengths differ by at most one. At global
depth $k$ the nonempty intervals have lengths $h_k$ or $h_k+1$, until they
are singletons, where $h_k=\max\{1,\lfloor n/2^k\rfloor\}$.

Consider a proper descendant $I$ of an ancestor $Q$, at global depth $k$,
and set $h_0=h_k$. Every side of $I$ belongs to $[h_0,2h_0]$.
Every positive clearance from $I$ to the boundary of $Q$ is a union of
depth-$k$ intervals, so is at least $h_0$. This uses the global $h_0$,
not the smallest side of this particular $I$: at odd sizes that substitution
can be false. If every side of $Q$ is at least two, a proper descendant
is shorter in every direction, so each axis has at least one positive
clearance. If the smallest side of $Q$ is one, every side is one or two;
each proper descendant is a corner singleton. That small case is treated
separately below.

For the large case put $\Omega=Q\setminus I$. Expand $I$ by exactly $h_0$
on every face with positive clearance, leaving a boundary-touching face
unchanged, and call the resulting rectangle $B$. Its side lengths $\ell_i$
lie in $[2h_0,4h_0]$. The collar $A=B\setminus I$ is covered by the actual
positive-clearance face slabs: normal width $h_0$ and all other sides those
of $B$. There are between $D$ and $J=2D$ slabs, each contained in $\Omega$.
For slabs $S_i,S_j$ with distinct normals, direct product counting gives
\begin{equation}\label{eq:balanced223-overlap}
 \frac{|S_i|}{|S_i\cap S_j|}=\frac{\ell_j}{h_0}\le4,
 \qquad \frac{|S_i|}{|S_j|}=\frac{\ell_j}{\ell_i}\le2.
\end{equation}
Choose the first slab as root; attach every distinct-normal slab directly
to it, and its opposite slab, if present, through the first slab with a
different normal. This tree has depth at most two.

\subsection{An unweighted BV extension, with every hole-crossing edge}
For a finite product rectangle $S$ of longest side at most $4h_0$, write
$V_1(S)=\sum_S|f-\operatorname{mean}_S f|$ and let $\mathrm{TV}_S(f)$
be the sum of $|df|$ over its actual internal edges. Then
\begin{equation}\label{eq:balanced223-L1}
 V_1(S)\le2h_0\mathrm{TV}_S(f).
\end{equation}
Indeed on an interval with $m$ sites, averaging pair paths bounds the
coefficient of its $r$th edge by $2r(m-r)/m\le m/2$.
Commuting coordinate-average projections contract $L^1$; telescoping their
product and applying this interval inequality fibrewise proves
\eqref{eq:balanced223-L1}, without an additional factor $D$.

Let $m_i$ be the mean on $S_i$ and let $c_0$ be the root mean.
Averaging on the actual intersection and using
\eqref{eq:balanced223-overlap} gives
$|S_i||m_i-m_j|\le4\{V_1(S_i)+V_1(S_j)\}$.
For a depth-two path $i,p,r$, the second edge costs the volume ratio
$|S_i|/|S_p|\le2$, so
\[
 |S_i||m_i-m_r|
 \le4V_1(S_i)+12V_1(S_p)+8V_1(S_r)
 \le12\sum_j V_1(S_j).
\]
Direct paths satisfy the same safe upper bound. Covering $A$ and counting
each actual edge at most $J$ times yields, with $P_D=2J(1+12J)$,
\begin{equation}\label{eq:balanced223-collar}
 \sum_A|f-c_0|\le(1+12J)\sum_jV_1(S_j)
 \le P_Dh_0\mathrm{TV}_A(f).
\end{equation}

At each positive-clearance hole face take the adjacent outside layer
of width $h_0$, with tangential sides those of $I$. It is contained in
$A$. On a normal fibre its nearest value satisfies
\[
 |f(\mathrm{first})-c_0|
 \le h_0^{-1}\sum_{\mathrm{fibre}}|f-c_0|
       +\sum_{\mathrm{fibre}}|df|.
\]
Sum over all actual faces and use \eqref{eq:balanced223-collar}.
The sum of hole-crossing jumps to the constant $c_0$ is at most
$J(P_D+1)\mathrm{TV}_A(f)$. Fill $I$ by $c_0$ and leave $f$ unchanged
on $\Omega$, obtaining $\widetilde f$ on $Q$. Internal hole differences
vanish; the added variation consists precisely of these crossing faces.
Therefore
\begin{equation}\label{eq:balanced223-BV}
 \mathrm{TV}_Q(\widetilde f)
 \le C_{\rm ext}\mathrm{TV}_\Omega(f).
\end{equation}
Boundary-touching faces create no exterior edge; at $h_0=1$ the fibre
gradient sum is zero. No coefficient is extended, reflected, averaged or
reassigned. In particular \eqref{eq:balanced223-BV} makes no weighted
quadratic trace-extension claim.

\subsection{Actual critical Sobolev and weighted gap on a ring}
Here are the finite details needed to pass from BV to the original energy.
For $T\subset\{0,\ldots,m-1\}^D$ with $|T|\le m^D/2$, the occupied-fibre
argument and the discrete projection inequality give
\[
 |\partial T|\ge\sum_i|\pi_iT|-D|T|/m
 \ge D|T|^{(D-1)/D}-D|T|/m
 \ge c_{\rm iso}|T|^{(D-1)/D},\quad
 c_{\rm iso}=D(1-2^{-1/D})\ge\tfrac12.
\]
Each occupied nonfull fibre supplies an edge, while full fibres number at
most $|T|/m$. The projection inequality
$|T|^{D-1}\le\prod_i|\pi_iT|$ is the finite Loomis--Whitney inequality;
it follows from iterated H\"older. A singleton cube has only zero variance.
Median parts, finite coarea and Minkowski consequently imply
\[
 \inf_z\|u-z\|_{D/(D-1)}\le(2/c_{\rm iso})\mathrm{TV}(u)
\]
on the cube. On a rectangle of aspect at most two, fold a cube of its
longest side onto it: in each coordinate first visit the interval in
order, repeat its final endpoint and reflect until the prescribed length.
Every vertex and edge has at most two preimages per coordinate; the
initial product block maps identically onto the rectangle. Restriction
and the variation multiplicity at most $2^D$ give the safe constant
$C_{\rm full}=2^{D+1}/c_{\rm iso}$ for that rectangle, still unweighted.

Apply this bound and \eqref{eq:balanced223-BV} to an indicator on $\Omega$.
For every real $z$,
$\|1_T-z\|_{s,\Omega}\ge\frac12\min\{|T|,|\Omega\setminus T|\}^{1/s}$,
where $s=D/(D-1)$. Thus $\Omega$ has relative edge isoperimetric constant
\[
 c_\Omega=c_{\rm iso}/(2^{D+2}C_{\rm ext}).
\]
Coarea gives $\|g\|_s\le c_\Omega^{-1}\mathrm{TV}_\Omega(g)$ for
nonnegative $g$ supported on at most half the vertices. With a median $m_f$
of $f$, take $w=(f-m_f)_+$ or $(m_f-f)_+$ and put
\[
 \alpha=2(D-1)/D,\qquad q=2D/(D+2),\qquad q'=2D/(D-2).
\]
Then $\alpha s=2$ and $(\alpha-1)q'=2$. The power-difference inequality,
H\"older and degree at most $2D$ yield
\[
 \mathrm{TV}(w^\alpha)
 \le2\alpha D^{1/q'}\|dw\|_q\|w\|_2^{\alpha-1}.
\]
In fact the $q'$th power of the endpoint factor sums to at most
$2^{q'}D\sum_xw(x)^2$. Divide only when $w\ne0$. The positive and
negative truncated $q$-edge energies sum to at most the original one;
since $2/q\ge1$, their squared bounds combine to give
\begin{equation}\label{eq:balanced223-q}
 \sum_\Omega|f-\operatorname{mean}_\Omega f|^2
 \le K_\Omega^2\Bigl(\sum_{\Omega\text{ edges}}|df|^q\Bigr)^{2/q},
 \qquad K_\Omega=2\alpha D^{1/q'}/c_\Omega.
\end{equation}
Only now use weights, on the original retained positive-start edges:
\[
 \Bigl(\sum |df|^q\Bigr)^{2/q}
 \le\Bigl(\sum_{\Omega\text{ edges}}a(\mathrm{start})^{-D/2}\Bigr)^{2/D}
       \E_a(f;\Omega)
 \le\{DS(\Omega)\}^{2/D}\E_a(f;\Omega).
\]
Here $q/(2-q)=D/2$, and at most $D$ retained edges start at one vertex.
Since $\alpha\le2$, $c_{\rm iso}\ge1/2$ and $2/q'+2/D=1$,
$K_\Omega^2D^{2/D}\le2^{2D+10}DC_{\rm ext}^{\,2}\le G_D$.
Full balanced rectangles obey the same safe bound without removing a hole.

For the small nodes with a singleton direction, $Q$, or $Q$ minus a corner,
is a connected graph with at most $b$ vertices whenever nonempty. Pair paths
of length at most $b-1$ give variance at most $b^2\E_1$.
Each internal edge satisfies
$a(\mathrm{start})^{-D/2}\le S(\Omega)$, so
$\E_1\le S(\Omega)^{2/D}\E_a$. As $G_D\ge b^2$, this proves in every case
\begin{equation}\label{eq:balanced223-gap}
 \sum_\Omega|f-\operatorname{mean}_\Omega f|^2
 \le G_DS(\Omega)^{2/D}\E_a(f;\Omega).
\end{equation}
A singleton has zero variance. Empty groups are never output.

\subsection{Coefficient-only heavy-chain compression and complete count}
Give each vertex mass $a(x)^{-P}$ and call a tree node heavy if its mass is
strictly greater than $t$. If the root is light, output it. Otherwise let
$L$ be the number of terminal heavy nodes and $B_h$ the number with at
least two heavy children. Terminal heavy nodes are disjoint, so
$L\le S(X)/t$. Finite tree counting gives $B_h\le L-1$.
After contracting unary paths, there are at most $B_h+L-1$ skeleton edges
and possibly a root-leading chain; hence at most $B_h+L\le2L-1$ maximal
unary chains. The bound is on chains, not their unrestricted total length.

At a unary node $Q\to Q'$ the increment $Q\setminus Q'$ is the union of
at most $b-1$ light children, of mass at most $(b-1)t$. Traverse a maximal
chain from ancestor toward descendant, closing a group when the cumulative
mass first reaches $t$. A closed group is a balanced ancestor ring
$Q_i\setminus Q_j$ of mass in $[t,bt)$; a residual is output only if nonempty,
and has mass less than $t$. Across all chains the disjoint closed groups
number at most $S(X)/t$, while residuals number at most $2L-1$.
At branching or terminal heavy internal nodes output each light child;
at a heavy singleton terminal output the singleton. Such cubes/singletons
number at most $b(B_h+L)\le b(2L-1)$.

These groups partition $X$: a unary ring excludes the endpoint heavy
subtree, which is then handled at the endpoint; branching processing
recurses into every heavy child and outputs all remaining children.
Thus their total number is at most
\[
 S(X)/t+(2L-1)+b(2L-1)\le1+(2b+3)S(X)/t.
\]
Each nonsingleton group has mass at most
$bt=(cG_D)^{-D/2}$, so \eqref{eq:balanced223-gap} gives residual gap at least
$c$. Heavy singleton mass may exceed this threshold, but its variance is
zero and its mean captures its entire norm.

\subsection{Rectangular rows while retaining intra-ring interfaces}
Partition $Q\setminus I$ by the first coordinate outside $I$: preceding
coordinates belong to $I$, the current coordinate is below or above $I$,
and later coordinates fill $Q$. This gives at most $2D$ nonempty integer
rectangles. Cubes and singletons require one rectangle. With
$\widehat q_A(f)=|A|^{-1/2}\sum_A f$, Cauchy gives
$\widehat q_\Omega(f)^2\le\sum_{B\subset\Omega}\widehat q_B(f)^2$.
The actual regional gap therefore implies
\[
 \E_a(f;\Omega)+c\sum_{B\subset\Omega}\widehat q_B(f)^2\ge c\|f\|_\Omega^2.
\]
Crucially $\E_a(f;\Omega)$ retains all edges between rectangular pieces
of this same ring; no separate gap is asserted for each piece.
Distinct regions are disjoint, so summing discards only nonnegative edges
between different regions. The full energy dominates their internal sum.
The count of rectangles is at most $2D$ times the region count, proving
Theorem~\ref{thm:balanced223}. The retained intra-ring interfaces are part of the regional energy throughout.


\begin{thebibliography}{9}
\bibitem{frank} R.~L.~Frank,
\emph{Cwikel's theorem and the CLR inequality}, arXiv:1206.3325v1 (2012),
Theorems~2.1 and~3.1. \url{https://arxiv.org/abs/1206.3325v1}.
\bibitem{brudnyi221} Yu.~Brudnyi,
\emph{Nonlinear Piecewise Polynomial Approximation and Multivariate $BV$
spaces of a Wiener--L.~Young Type.~I}, arXiv:1511.03971v1 (2015).
\url{https://arxiv.org/abs/1511.03971v1}.
\bibitem{cdpx223} A.~Cohen, R.~DeVore, P.~Petrushev and H.~Xu,
\emph{Nonlinear approximation and the space $BV(\R^2)$},
American Journal of Mathematics \textbf{121} (1999), 587--628.
Author PDF: \url{https://people.math.sc.edu/pencho/Publications/BV9707.pdf}.
\end{thebibliography}
\end{document}
